\section{El protocolo: rojo, azul y morado}\label{sec:protocolo}

Quien construye un detector y además lo evalúa tiene un incentivo permanente a ajustarlo hasta que
pase. El protocolo existe para que ese ajuste no sea posible, y para que quede evidencia de que no
ocurrió. Esta sección describe sus reglas, lo que se hizo en cada paso y lo que alguien de afuera
puede verificar.

\subsection{Las seis reglas}

Están escritas en \cod{docs/PLAN.md} desde antes de empezar. Las copiamos tal cual.

\begin{enumerate}
  \item «El rojo no mira el código del detector antes de H20. El agente del área Rojo no abre
    \cod{engine/} hasta entonces.»
  \item «El set de prueba lo genera el rojo, guarda su hash y queda sellado hasta H20. Sus familias no
    aparecen en entrenamiento.»
  \item «Revisiones de 15 minutos en H6, H12 y H17, solo sobre fallas de entrenamiento y validación.»
  \item «El azul puede corregir errores y ajustar parámetros en validación. No puede agregar reglas
    para un ataque concreto. Si lo hace, se anota.»
  \item «En H20 se congela. El rojo recibe la configuración completa y tiene 2 horas de ataque
    adaptativo, que se reporta aparte y no se usa para reajustar.»
  \item «Bitácora con hora y motivo de cada cambio al detector.»
\end{enumerate}

H20 es la hora 20 del evento, las 19:00 del sábado. El líder decidió adelantar el congelamiento a
las 09:40. Adelantarlo no relaja ninguna regla: el orden entre sellar, congelar y abrir es el mismo.

\subsection{Quién es cada equipo}

No son tres grupos de personas. Son tres papeles.

\begin{description}
  \item[Rojo.] Una sesión de agente aparte, que escribió los generadores de ataques y de sesiones
    normales difíciles. Su instrucción de trabajo incluye no abrir \cod{engine/}. Eligió las familias
    del set de prueba a partir de técnicas documentadas en fuentes públicas, no de lo que el detector
    pudiera ver.
  \item[Azul.] Los agentes que construyeron el detector, ajustaron el modelo y buscaron los
    parámetros, con entrenamiento y validación.
  \item[Morado.] Las reglas de arriba y los mecanismos que las hacen cumplir: el sello, el
    congelamiento con hash, el candado de la corrida única y la bitácora.
\end{description}

\begin{limite}
\textbf{Lo que no se puede comprobar desde fuera.} Que la sesión del rojo no abrió el detector es una
regla de trabajo. No hay un mecanismo técnico que lo demuestre. Lo que sí es comprobable es todo lo
demás: el orden de los hechos, los hashes y que el comando de evaluación no se puede repetir.
\end{limite}

\subsection{La bitácora}

\cod{docs/BITACORA_DETECTOR.md} lista cada cambio al detector con hora, qué cambió y por qué. Entre la
construcción inicial y el congelamiento hay nueve entradas. Tres importan para leer los resultados.

\begin{itemize}
  \item \textbf{04:05.} El período $T$ baja de 200 a 150 y se agregan la invariante, el tope de tres
    rotaciones y el modo escalado. Motivo: un goteo de 7 caracteres sacaba hasta 112 bits.
  \item \textbf{06:30.} Se agrega la cuenta compartida. Motivo: con cinco agentes salían los 160 bits.
    Es una capacidad general.
  \item \textbf{07:30.} El contador pasa a leer base64 corto. La bitácora lo marca como el único cambio
    que \emph{sí} responde a un ataque concreto: una familia de entrenamiento y la demo interactiva.
    Es lo que la regla 4 pide anotar, y está anotado.
\end{itemize}

\subsection{El ajuste y la búsqueda}\label{sec:congelamiento}

El detector tiene dos partes que salen de datos: el modelo de evidencia y cuatro umbrales.

\paragraph{El modelo.} Se ajusta solo con entrenamiento (\cref{sec:modelo-ajustado}).

\paragraph{Los umbrales.} Se eligen solo con validación, con \cod{python -m eval.search}. La búsqueda
recorre una grilla de 750 combinaciones:
\begin{gather*}
  b \in \{4;\ 6;\ 8;\ 9{,}2;\ 11;\ 13\}, \qquad
  \brot \in \{2;\ 3;\ 4;\ 5;\ 6\}, \\
  \lrot \in \{0{,}1;\ 0{,}2;\ 0{,}3;\ 0{,}4;\ 0{,}5\}, \qquad
  T \in \{25;\ 50;\ 75;\ 100;\ 150\} .
\end{gather*}
Para cada una reproduce el set de validación (50 sesiones normales y 80 ataques) y calcula cuatro
métricas: falsas alarmas, rotaciones en sesiones normales por cada \num{1000} llamadas, ataques
frenados y mediana de bits filtrados.

La elección sigue un criterio de Neyman--Pearson: primero se fija el costo tolerable y después se
minimiza el daño.

\begin{enumerate}
  \item Quedan las combinaciones con falsas alarmas $\leq \num{0.001}$ y rotaciones en sesiones
    normales $\leq 5$ por cada \num{1000} llamadas. Cumplen 450 de las 750.
  \item Entre esas, se elige la de menor mediana de bits filtrados y, en caso de empate, la de más
    ataques frenados.
\end{enumerate}

El resultado fue $b = 4$, $\brot = 2$, $\lrot = \num{0.1}$, $T = 75$: cero falsas alarmas, 0,47
rotaciones por cada \num{1000} llamadas normales, todos los ataques de validación frenados y una
mediana de 16 bits.

\paragraph{Qué tan determinados quedaron.} Las \cref{tab:busqueda-rot,tab:busqueda-bits} resumen las
750 filas de \cod{models/search/rows.json}, y muestran algo que conviene decir con claridad.

\begin{table}[ht]
\caption{Rotaciones en sesiones normales por cada \num{1000} llamadas, en validación, según $T$ y
$\brot$ (rango sobre los demás parámetros). La meta era 5 o menos.}\label{tab:busqueda-rot}
\small
\begin{tabular}{@{}rccccc@{}}
\toprule
 & \multicolumn{5}{c}{$\brot$} \\
\cmidrule(l){2-6}
$T$ & 2 & 3 & 4 & 5 & 6 \\
\midrule
\tabladatos{tab_busqueda_rot}
\bottomrule
\end{tabular}
\end{table}

\begin{table}[ht]
\caption{Mediana de bits filtrados en los ataques de validación, según $\lrot$ y $\brot$ (rango sobre
los demás parámetros), y porcentaje de ataques frenados según $\brot$.}\label{tab:busqueda-bits}
\small
\begin{tabular}{@{}lccccc@{}}
\toprule
 & \multicolumn{5}{c}{$\brot$} \\
\cmidrule(l){2-6}
$\lrot$ & 2 & 3 & 4 & 5 & 6 \\
\midrule
\tabladatos{tab_busqueda_bits}
\bottomrule
\end{tabular}
\end{table}

\begin{itemize}
  \item \textbf{El período decide el costo.} Con $T$ de 25 o 50 hay demasiadas rotaciones periódicas
    en sesiones normales; desde 75 casi desaparecen.
  \item \textbf{$\brot$ tiene dos regímenes.} Con $\brot \leq 3$, una sola llamada marcada rota, y se
    frenan todos los ataques de validación con 16 bits de mediana. Con $\brot \geq 4$ hacen falta dos
    llamadas marcadas seguidas, y algunos ataques se escapan. Es la consecuencia directa de un modelo
    cuyo incremento vale $+3$ o $-3$.
  \item \textbf{$b$ no cambia nada.} Ninguna de las cuatro métricas se mueve al variar el umbral de
    bloqueo por evidencia: la regla 3 no se disparó nunca en validación.
  \item \textbf{La combinación elegida no es única.} Hay 180 combinaciones que empatan con ella en
    los dos criterios de la elección (16 bits de mediana y todos los ataques frenados). El criterio se
    queda con la primera en el orden de la grilla, y por eso salieron los valores más bajos de $b$,
    $\brot$ y $\lrot$ y el menor $T$ que cumple la meta.
  \item \textbf{Y no es la mejor en validación.} De esas 180, 144 no rotaron nunca en una sesión
    normal; la elegida rotó una vez (0,47 por cada \num{1000} llamadas). Las 144 son exactamente la
    frontera de Pareto que guarda \cod{models/search/pareto.json}, y la configuración congelada no
    está en ella. La diferencia es $\lrot$: \num{0.1} es el único valor de la grilla con el que hubo
    una rotación de más, porque con 16 bits de umbral basta una coincidencia de 4 caracteres por azar
    (\cref{sec:contador-limites}). El criterio de elección no mira las rotaciones una vez que están
    bajo la meta, y por eso no lo distinguió.
\end{itemize}

\begin{limite}
\textbf{Cómo leer la elección.} Validación tiene las mismas cinco familias que entrenamiento,
generadas por el mismo programa con otra semilla. Sus cifras son una cota optimista y no separan
bien las configuraciones. De los cuatro umbrales, los datos determinan dos decisiones gruesas
($T \geq 75$ y $\brot \leq 3$). El valor de $b$ quedó puesto por el orden de la grilla, y el de
$\lrot$ también, aunque un valor mayor habría rotado menos en validación. Esto se vio al escribir
este informe, con el detector ya congelado y evaluado: no se cambió nada, y no se midió otra
configuración sobre el set de prueba, porque eso sería ajustar mirando la prueba. La cifra que vale
es la del set de prueba, con la configuración que se congeló.
\end{limite}

\subsection{El congelamiento}

Congelar es escribir \cod{models/search/frozen.json}, que contiene:

\begin{itemize}
  \item los cuatro umbrales;
  \item el sha256 del archivo del modelo (\cod{993976e7}…);
  \item un sha256 de las carpetas de datos usadas;
  \item la fecha y la hora;
  \item y un \emph{hash} de todo lo anterior: \cod{cf56ce20bb718b1c3a3297cf4984ae76c839779acb01a6d47df687fc08bc50ba}.
\end{itemize}

El detector evaluado se construye siempre con \cod{make_frozen}, que antes de nada vuelve a calcular el
sha256 del modelo en disco y se niega a continuar si no es el congelado. Cada resultado guarda en su
\cod{RUN.json} el hash con que corrió: los cinco conjuntos de resultados del repositorio llevan el
mismo.

\subsection{La corrida única}

La evaluación sobre el set de prueba es un solo comando,
\cod{python -m eval.final.run --split test --out results/test}, con un candado: lo primero que hace es
escribir \cod{results/test/RUN.json} con el estado «en curso», y si ese archivo ya existe se niega a
correr. No tiene opción para forzarlo. Como el archivo se escribe al empezar, ni siquiera una corrida
interrumpida se puede repetir sin borrarlo a mano, a sabiendas.

La corrida empezó a las 09:40:26 y terminó a las 09:41:19, en el commit \cod{cd14e9e}: 360 sesiones,
300 unidades y nueve sistemas en 52,6 segundos.

\begin{table}[ht]
\caption{La secuencia verificable del sábado 21 de noviembre. Los commits están en el historial del
repositorio.}\label{tab:secuencia}
\small
\begin{tabularx}{\textwidth}{@{}P{1.4cm}LP{4.6cm}@{}}
\toprule
Hora & Hecho & Cómo se comprueba \\
\midrule
07:10 & El rojo publica el sello del set de prueba & \cod{redteam/SELLO.md}; commit \cod{77457d9} \\
09:36 & El azul congela la configuración & \cod{frozen.json}; commit \cod{3f26307}, etiqueta
  \cod{congelamiento} \\
09:40 & Se abre el set y el sello coincide & El comando de \cod{SELLO.md} \\
09:40 & La evaluación corre una sola vez & \cod{results/test/RUN.json}; commit \cod{cd14e9e} \\
09:53 & Análisis del atacante que se adapta & \cod{results/adaptive/RUN.json} \\
14:07 & Análisis con las marcas del proxy & \cod{results/proxy_flags/RUN.json}; commit \cod{8f1ec16} \\
15:32 & Análisis de los fragmentos cortos & \cod{results/short_fragments/RUN.json} \\
\bottomrule
\end{tabularx}
\end{table}

Los tres análisis posteriores usan el detector congelado, se reportan aparte y no cambiaron nada.

\subsection{Unidades y sistemas}

\paragraph{Unidades.} La unidad de evaluación no es siempre la sesión. En la familia A7 un ataque son
tres sesiones que comparten una clave; contarlas por separado sería contar tres ataques fallidos donde
hubo uno. Las sesiones cuyas etiquetas comparten un grupo forman una \emph{unidad}: se reproducen en
orden y, antes de medir, se funden en una sola sesión lógica. Toda otra sesión es una unidad por sí
misma.

\paragraph{Sistemas.} Además de Alicanto se miden tres ablaciones suyas y seis líneas de base, todas
sobre las mismas unidades y con la misma métrica (\cref{tab:que-sistemas}).

\begin{table}[ht]
\caption{Los sistemas que se comparan. Las líneas de base están en \cod{eval/baselines/}.}\label{tab:que-sistemas}
\small
\begin{tabularx}{\textwidth}{@{}P{4.6cm}L@{}}
\toprule
Sistema & Qué es \\
\midrule
Alicanto (congelado) & El detector congelado. Por unidad, una cuenta compartida nueva y un detector
  por sesión \\
Con una cuenta por sesión & El mismo detector, sin cuenta compartida \\
Sin evidencia & Los mismos parámetros y la misma cuenta, sin el modelo de evidencia \\
Sin rotación & Ninguna regla rota; los bloqueos siguen \\
\midrule
Regla simple & «La sesión leyó una clave (su valor completo en un resultado) y esta llamada saca
  datos hacia un destino no confiable»: rota lo leído; nunca bloquea \\
Lista de destinos permitidos & Bloquea todo sumidero hacia un \esp{host} fuera de una lista de cinco \\
Escáner de secretos & La biblioteca \cod{detect-secrets}, con sus detectores por defecto, sobre los
  argumentos de cada sumidero. No conoce las claves \\
Lista más escáner & Bloquea si cualquiera de las dos anteriores bloquea \\
Presupuesto de bits (réplica) & El mismo contador de Alicanto, por sesión y sin rotación: al llegar a
  96 bits bloquea esa llamada y todos los sumideros siguientes \\
Reensamblado (réplica) & Junta lo que sale por cada campo a lo largo de la sesión, sin decodificar, y
  bloquea al ver 16 caracteres seguidos de una clave \\
\bottomrule
\end{tabularx}
\end{table}

Las dos réplicas son versiones simplificadas que hicimos a partir de documentación pública de
herramientas que ya juntan fragmentos entre llamadas. No son esos productos ni su código, y sus cifras
no dicen nada de ellos. Están para tener, al menos, una referencia de las dos ideas más cercanas a la
nuestra.

\subsection{Las métricas}\label{sec:metricas}

Todas se calculan desde el registro de eventos y las etiquetas (\cod{eval/metrics/metrics.py}).

\begin{description}
  \item[Llamada expuesta.] Una llamada es expuesta si saca datos, o si ocurre desde que la sesión leyó
    por primera vez una clave o contenido no confiable. Son las llamadas en que un bloqueo erróneo es
    posible; las demás no cuentan en el denominador.
  \item[Falsa alarma.] Un \emph{bloqueo} de una llamada expuesta en una sesión normal. La tasa es el
    número de falsas alarmas dividido por el de llamadas expuestas de las sesiones normales. Una
    rotación no es una falsa alarma: no interrumpe nada.
  \item[Rotaciones en sesiones normales.] Decisiones con rotación en sesiones normales, por cada
    \num{1000} llamadas normales. Es el costo propio de Alicanto y se informa aparte.
  \item[Bits que salieron.] Para cada clave de la unidad se acumulan las posiciones que las etiquetas
    dicen que salieron, \emph{solo} en llamadas que el detector dejó pasar. Cuando el detector rota la
    clave, el acumulado vuelve a cero. El valor que se informa es el máximo del acumulado a lo largo de
    la unidad, por 4 bits.
  \item[Ataque frenado.] Para toda clave de la unidad, ese máximo no pasa de $H - \beta = 96$ bits.
\end{description}

\paragraph{Por contenido.} Hay un punto fino en cómo se cuenta después de una rotación. La grabación
sigue enviando pedazos de la clave original, que ya no es la vigente. La regla principal es contar
\emph{por contenido}: un pedazo suma solo si es de la clave vigente, así que después de la primera
rotación los envíos de la grabación no suman. Es lo correcto para un atacante que no se adapta, y hay
que leerlo así: mide cuánto de una clave útil salió.

Se informa además una \emph{cota pesimista}: tras rotar, los envíos siguientes cuentan contra la clave
nueva si la sesión volvió a leer el archivo de claves. Es el caso en que el atacante podría haber
obtenido el valor nuevo. En el set de prueba las dos cuentas dan lo mismo.

\subsection{Las pruebas estadísticas}

La evaluación final incluye cuatro pruebas, corregidas por Holm~\cite{holm}
(\cod{eval/final/tests_h.py}).

\begin{table}[ht]
\caption{Las pruebas de la evaluación final. T5 y T6, previstas en el diseño, no se calcularon.}\label{tab:pruebas-def}
\small
\begin{tabularx}{\textwidth}{@{}P{0.8cm}LL@{}}
\toprule
 & Qué afirma & Cómo se prueba \\
\midrule
T1 & La tasa de falsas alarmas es menor que \num{0.001} & Cota superior unilateral al
  \qty{95}{\percent} de Clopper--Pearson~\cite{clopper}. Con cero alarmas hacen falta unas
  \num{3000} llamadas expuestas para pasar \\
T2 & Se frena más del \qty{80}{\percent} de los ataques & Cota inferior unilateral al
  \qty{95}{\percent} de Clopper--Pearson \\
T3 & Alicanto frena ataques que la línea de base no frena & Prueba exacta de McNemar~\cite{mcnemar}
  sobre los pares de unidades, contra «lista más escáner» \\
T4 & La rotación reduce el éxito del ataque & Prueba exacta de Fisher sobre el éxito con y sin
  rotación; Mann--Whitney sobre los bits \\
\bottomrule
\end{tabularx}
\end{table}

\begin{limite}
\textbf{Los valores p exageran.} Las pruebas tratan las 180 unidades como independientes, y son 30
réplicas de seis plantillas. Un intervalo de confianza estrecho alrededor de «120 de 180» no dice
que el sistema frenaría esa fracción de ataques nuevos: dice que frena esas familias de forma
consistente. El resultado se lee por familias.
\end{limite}
